% -*- Mode: LaTeX; coding: utf-8; indent-tabs-mode: nil; tab-width: 2 -*-

%% This code is part of the Introduction to Tests talk.
% See https://github.com/gonzalo-bulnes/introduction-to-tests for the last version.
%
% You should'nt compile this file directly, as it will be included by
% the other .tex files.
%%

\begin{frame}[t,plain]
  \maketitle
\end{frame}

\begin{frame}
  \tableofcontents
\end{frame}

% Licencias de software libre
\section{\dots}

  \begin{frame}
    \tableofcontents[currentsection]
  \end{frame}

%
\subsection{Definiciones}

  \begin{frame}{\insertsection}
    % TODO
    \dots
  \end{frame}


% Conclusión
\section*{Conclusión}

  \begin{frame}{\insertsection}
    % TODO
  \end{frame}

\appendix

% Un poco de lectura
\section{Un poco de lectura}

%
\subsection{Referencias}

  \begin{frame}{\insertsubsection}
    \bibliography{../bibtex/biblio}{}
    \bibliographystyle{plain}
    % Learn Python!
    %\nocite{pilgrim2004dive}
    %\nocite{pink2011drive}
  \end{frame}

%
\subsection{Este documento es libre}

  % Sources of images from other authors:
  % (Propriety of their authors!)
  % http://www.acid.cl (logo created by Francisco Vásquez <fvasquez@kolorsheep.com>)

  \begin{frame}{\insertsubsection}

    \footnotesize
      \begin{columns}
        \column{.6\textwidth}
          Este \href{https://github.com/gonzalo-bulnes/introduction-to-tests}{documento}, sus \href{hhttps://github.com/gonzalo-bulnes/introduction-to-tests}{fuentes}, y el \href{https://github.com/gonzalo-bulnes/santiago-beamer-theme}{tema para Beamer llamado Santiago} están disponibles en:\\~\\
          \url{https://github.com/gonzalo-bulnes}
        \column{.3\textwidth}
        \hfill\includegraphics[width=\textwidth]{../../src/images/sources-talk.png}
      \end{columns}
    \normalsize

    \scriptsize{Menos mención contraria:}\\
    \noindent\scriptsize{Copyright~\copyright~2013~Gonzalo~Bulnes~Guilpain}\\
    \noindent\scriptsize{Copyright~\copyright~2013~Acid~Ltda}

    \scriptsize{Permission is granted to copy, distribute and/or modify this document
    under the terms of the GNU Free Documentation License, Version 1.3
    or any later version published by the Free Software Foundation;
    with no Invariant Sections, no Front-Cover Texts, and no Back-Cover Texts.
    A copy of the license can be found at \url{http://www.gnu.org/copyleft/fdl.html}.}
  \end{frame}

%
\subsection{Agradecimientos}

  \begin{frame}
    \vspace{0mm}
    \begin{center}
      Gracias por su atención.
    \end{center}
  \end{frame}
% EOF
